\documentclass[../../../../main.tex]{subfiles}
\begin{document}

\begin{flushright}
\footnotesize\textit{【自動文字起こし・要確認】}
\end{flushright}

[解] $C_1, C_2$ の中心を $O_1, O_2$ とする。以下 $C = \cos\theta, S = \sin\theta$ とする。
\[
\begin{cases}
|\text{P}O_1|^2 = (C+(1-a))^2 + S^2 = (a-1)^2+1+2(1-a)C \\
|\text{P}O_2|^2 = (C-(1-b))^2 + S^2 = (b-1)^2+1-2(1-b)C
\end{cases}
\]
だから $\Delta \text{PQ}O_1, \Delta \text{PR}O_2$ にピタゴラスの定理を用いて、
\begin{align*}
|\text{PQ}|^2 &= |\text{P}O_1|^2 - a^2 = 2(1-a)(1+C) \\
|\text{PR}|^2 &= |\text{P}O_2|^2 - b^2 = 2(1-b)(1-C)
\end{align*}
だから、$|\text{PQ}|, |\text{PR}| \ge 0$ より、
\[
|\text{PQ}| = \sqrt{2(1-a)(1+C)}, \quad |\text{PR}| = \sqrt{2(1-b)(1-C)}
\]
である。コーシー・シュワルツより
\begin{align*}
|\text{PQ}|+|\text{PR}| &\le \sqrt{(1+1)\{2(1-a)(1+C) + 2(1-b)(1-C)\}} \\
&= 2\sqrt{2-a-b} \quad \cdots \star
\end{align*}
で、等号成立は、$0 \le a, b < 1$ より、（$\because C_1, C_2$ が $C$ に含まれる）
\begin{align*}
&\sqrt{2(1-a)(1+C)} = \sqrt{2(1-b)(1-C)} \\
\Leftrightarrow\ & (1-a)(1+C) = (1-b)(1-C) \\
\Leftrightarrow\ & (2-a-b)C = a-b \quad \cdots \text{①}
\end{align*}
であり、$2-a-b \ge a-b \Leftrightarrow 1 \ge a$ より、①をみたす $C$ （$-1 \le C \le 1$）が必ず存在する。以上から、$\star$より、
\[
\max\{|\text{PQ}|+|\text{PR}|\} = 2\sqrt{2-a-b}
\]

\end{document}
